% Source printed pp. 185-195 / PDF pp. 193-203.
\section{粗模空间与 Keel--Mori 定理}\label{src:5.4}
\begin{epigraph}
至于 $M_g$, 几乎毫无疑问, 可以赋予它代数簇的结构.
\epigraphauthor{André Weil~\cite[第 $383$ 页]{Wei58}}
\end{epigraph}
我们将证明, Noether 概形上的每个分离有限型 Deligne--Mumford 叠 $\mathcal X$ 都有分离的粗模空间 $\pi:\mathcal X\to X$, 称为 Keel--Mori 定理 (\bookref{5.4.6}). 可将此定理视为消去 Deligne--Mumford 叠的叠性的一种方式: 以牺牲 $\mathcal X$ 的泛性质 (例如万有族的存在性) 为代价, 可以用代数空间替换 $\mathcal X$, 而不改变底层拓扑空间. 引入稳定曲线的 Deligne--Mumford 叠 $\overline{\mathcal M}_g$ (定理~\bookref{6.4.16}) 并证明它分离 (命题~\bookref{6.5.20}) 后, Keel--Mori 定理蕴含存在粗模空间 $\overline{\mathcal M}_g\to\overline M_g$, 其中 $\overline M_g$ 是分离代数空间. 此外, 半稳定约化 (\bookref{6.5.9}), 即 $\overline{\mathcal M}_g$ 的固有性, 进一步蕴含 $\overline M_g$ 固有. 在~\S\bookref{6.8} 和~\S\bookref{6.9} 中, 我们将给出 $\overline M_g$ 射影性的两个证明.

粗模空间已在定义~\bookref{5.2.1} 中引入. 若 $G$ 是作用于仿射概形 $\Spec A$ 的有限抽象群, 则 $[\Spec A/G]\to\Spec A^G$ 是粗模空间 (定理~\bookref{5.2.3}), 换言之 $\Spec A\to\Spec A^G$ 是几何商. 为证明 Keel--Mori 定理 (\bookref{5.4.6}), 将应用局部结构定理 (\bookref{5.3.1}) 构造平展邻域 $[\Spec A_i/G]\to\mathcal X$, 再在平展拓扑中粘合几何商 $\Spec A_i^G$, 从而构造 $\mathcal X$ 的粗模空间.

\subsection{将平展态射下降到几何商}\label{src:5.4.1-subsection}
为在平展拓扑中粘合几何商, 需要一个条件, 保证平展 $G$-等变态射 $\Spec A\to\Spec B$ 在几何商上诱导平展态射 $\Spec A^G\to\Spec B^G$. 以下命题可视为 Luna 基本引理 (\bookref{7.5.30}) 的初步版本.
\begin{proposition}\label{src:5.4.1}
设 $G$ 是有限抽象群, $f:\Spec A\to\Spec B$ 是 Noether 环 $R$ 上有限型仿射概形的 $G$-等变态射. 设 $x\in\Spec A$ 是闭点. 假设
\begin{enumerate}[(a)]
\item $f$ 在 $x$ 处平展, 且
\item 所诱导的稳定子群概形的映射 $G_x\to G_{f(x)}$ 是双射.
\end{enumerate}
则 $x$ 的像在 $\Spec A^G$ 中有仿射开邻域 $W$, 使得 $W\to\Spec A^G\to\Spec B^G$ 平展, 且 $\pi_A^{-1}(W)\cong W\times_{\Spec B^G}[\Spec B/G]$, 其中 $\pi_A:[\Spec A/G]\to\Spec A^G$.
\end{proposition}
\begin{remark}\label{src:5.4.2}
换言之, 将 $\Spec A^G$ 换成 $\pi_A(x)$ 的仿射邻域 $W$, 并将 $\Spec A$ 换成 $\pi_A^{-1}(W)$ 后, 可使图
\begin{equation}\label{src:5.4.3}
\begin{tikzcd}
{[\Spec A/G]}\arrow[r,"f"]\arrow[d,"\pi_A"']&{[\Spec B/G]}\arrow[d,"\pi_B"]\\
\Spec A^G\arrow[r]\arrow[ur,phantom,"\square" description]&\Spec B^G
\end{tikzcd}
\end{equation}
为 Cartesian 图, 其中两个水平映射均平展.
\end{remark}
\begin{proof}[命题~\bookref{5.4.1} 的证明]
置 $y=f(x)$. 首先断言, 问题在 $\pi_B(y)\in\Spec B^G$ 附近是平展局部的. 事实上, 若 $Y'\to Y:=\Spec B^G$ 是 $\pi_B(y)$ 的仿射平展邻域, 用 $X'$, $\mathcal X'$ 和 $\mathcal Y'$ 分别表示 $X:=\Spec A^G$, $\mathcal X:=[\Spec A/G]$ 和 $\mathcal Y:=[\Spec B/G]$ 的基变换. 由引理~\bookref{5.2.7}, 有 $\mathcal Y'\cong[\Spec B'/G]$, $Y'\cong\Spec B'{}^G$, 对 $\mathcal X'$ 与 $X'$ 也类似. 若结果在此基变换后成立, 则存在包含 $\pi_A(x)$ 的某个原像的开邻域 $W'\subseteq X'$, 使得 $W'\hookrightarrow X'\to Y'$ 平展, 且 $W'$ 在 $\mathcal X'$ 中的原像同构于 $W'\times_{Y'}\mathcal Y'$. 取 $W$ 为 $W'$ 在 $X'\to\Spec A^G$ 下的像, 再应用平展下降, 即得所需断言.

现在断言, 由此可假设 $B^G$ 是严格 Hensel 环. 为此, 定义 $Y^{\mathrm{sh}}=\Spec\mathcal O^{\mathrm{sh}}_{Y,\pi_B(y)}$, 用 $X^{\mathrm{sh}}$, $\mathcal X^{\mathrm{sh}}$ 和 $\mathcal Y^{\mathrm{sh}}$ 表示 $X$, $\mathcal X$ 和 $\mathcal Y$ 沿 $Y^{\mathrm{sh}}\to Y$ 的基变换. 可将 $Y^{\mathrm{sh}}=\lim_\lambda Y_\lambda$ 写成 $y$ 的仿射平展邻域 $Y_\lambda\to Y$ 的极限, 用 $X_\lambda$, $\mathcal X_\lambda$ 和 $\mathcal Y_\lambda$ 表示 $X$, $\mathcal X$ 和 $\mathcal Y$ 沿 $Y_\lambda\to Y$ 的基变换. 假设 $U^{\mathrm{sh}}\subseteq X^{\mathrm{sh}}$ 是包含位于 $x$ 及 $Y^{\mathrm{sh}}$ 的闭点上方的唯一点的仿射开子概形, 使得 $U^{\mathrm{sh}}\to Y^{\mathrm{sh}}$ 平展, 且
\[
\pi_{\mathcal X^{\mathrm{sh}}}^{-1}(U^{\mathrm{sh}})\cong U^{\mathrm{sh}}\times_{Y^{\mathrm{sh}}}\mathcal Y^{\mathrm{sh}}.
\]
由命题~\bookref{B.3.3}, 态射 $U^{\mathrm{sh}}\to X^{\mathrm{sh}}$ 对某个 $\eta$ 下降为 $U_\eta\to X_\eta$. 对 $\lambda>\eta$ 置 $U_\lambda=U_\eta\times_{X_\eta}X_\lambda$. 由命题~\bookref{B.3.7}, 对 $\lambda\gg0$ 有
\begin{enumerate}[(a)]
\item $U_\lambda\to X_\lambda$ 是开浸入;
\item 复合 $U_\lambda\to X_\lambda\to Y_\lambda$ 平展; 且
\item $\pi_{\mathcal X_\lambda}^{-1}(U_\lambda)\cong U_\lambda\times_{Y_\lambda}\mathcal Y_\lambda$ (在 $\mathcal X$ 与 $\mathcal Y$ 的平展呈示上论证).
\end{enumerate}
最后, 由于 $B^G\to B$ 有限 (引理~\bookref{5.2.6}), $B=B_1\times\cdots\times B_r$ 是严格 Hensel 局部环的乘积 (命题~\bookref{B.5.10}). 如定理~\bookref{5.2.3} 的证明, 可将 $[\Spec B/G]$ 换成 $[\Spec B_1/G_y]$, 将 $[\Spec A/G]$ 换成 $[f^{-1}(\Spec B_1)/G]$, 从而假设 $G$ 固定 $x$ 与 $y$, 且在剩余域 $\kappa(x)=\kappa(y)$ 上作用平凡. 因此 $\Spec A\to\Spec B$ 有将 $y$ 送到 $x$ 的唯一截面 $s:\Spec B\to\Spec A$. 截面 $s$ 必为 $G$-不变的 (如定理~\bookref{5.2.3} 的证明). 于是 $s$ 下降为 $\Spec A^G\to\Spec B^G$ 的截面, 给出所需的开闭子概形 $W\subseteq\Spec A^G$.
\end{proof}
\begin{remark}\label{src:5.4.4}
以下从概念上解释为何应当期待所诱导的商映射平展. 为简单起见, 假设 $R=k$ 是代数闭域. 设 $\widehat A$ 与 $\widehat B$ 是 $x$ 与 $f(x)$ 处局部环的完备化. 稳定子 $G_x$ 与 $G_{f(x)}$ 分别作用于 $\Spec\widehat A$ 与 $\Spec\widehat B$, 映射 $\Spec\widehat A\to\Spec\widehat B$ 关于映射 $G_x\to G_{f(x)}$ 等变. $A^G$ 在 $x$ 的像处的完备化 $\widehat{A^G}$ 同构于 $\widehat A^{G_x}$ (习题~\bookref{5.2.15}), 类似地 $\widehat{B^G}=\widehat B^{G_{f(x)}}$. 由于 $f$ 在 $x$ 处平展, $\widehat B\to\widehat A$ 是同构, 又因 $G_x\to G_{f(x)}$ 是双射, 所诱导的映射 $\widehat{B^G}\to\widehat{A^G}$ 是同构, 这说明 $\Spec A^G\to\Spec B^G$ 在 $x$ 的像处平展.
\end{remark}
上述命题将在 Keel--Mori 定理 (\bookref{5.4.6}) 的证明中以下列形式应用.
\begin{corollary}\label{src:5.4.5}
设 $G$ 是有限抽象群, $f:\Spec A\to\Spec B$ 是 Noether 环 $R$ 上有限型仿射概形的 $G$-等变态射. 假设对每个闭点 $x\in\Spec A$,
\begin{enumerate}[(a)]
\item $f$ 在 $x$ 处平展, 且
\item 所诱导的稳定子群概形的映射 $G_x\to G_{f(x)}$ 是双射.
\end{enumerate}
则 $\Spec A^G\to\Spec B^G$ 平展, 且~\bookeqref{5.4.3} 是 Cartesian 图.
\end{corollary}

\subsection{Keel--Mori 定理}\label{src:5.4.2-subsection}
现在陈述并证明 Keel--Mori 定理.
\begin{theorem}[Keel--Mori 定理]\label{src:5.4.6}
设 $\mathcal X$ 是 Noether 代数空间 $S$ 上的分离有限型 Deligne--Mumford 叠. 则存在粗模空间 $\pi:\mathcal X\to X$, 满足 $\mathcal O_X=\pi_*\mathcal O_{\mathcal X}$, 且
\begin{enumerate}[(1)]
\item $X$ 在 $S$ 上分离且有限型;
\item $\pi$ 是固有万有同胚; 且
\item 对代数空间的每个平坦态射 $X'\to X$, 基变换 $\mathcal X\times_X X'\to X'$ 是粗模空间.
\end{enumerate}
此外, $X$ 固有, 当且仅当 $\mathcal X$ 固有.
\end{theorem}
此定理可推广到满足 $I_{\mathcal X}\to\mathcal X$ 有限的 Deligne--Mumford 叠 $\mathcal X$ (可能不分离), 只是其粗模空间可能不分离 (习题~\bookref{5.4.10}). 更为有力的是, Keel--Mori 定理~\cite{KM97} 在将 $\mathcal X$ 是 Deligne--Mumford 叠的假设换成 $\mathcal X$ 是具有有限惯性 $I_{\mathcal X}\to\mathcal X$ 的代数叠这一假设后仍成立; 见注记~\bookref{5.4.11}. 这要求所有稳定子有限, 但不一定约化. 此定理在不加任何 Noether 性或有限性条件时也成立; 见~\cite{Con05b}, \cite{Ryd13} 和~\cite[Tag 0DUK]{SP}.
\begin{remark}[历史注记]\label{src:5.4.7}
大约与~\cite{KM97} 发表同时, Kollár 证明了一个重要特例, 已足以用于许多应用: 若 $G$ 是优良概形 $S$ 上的仿射群概形, 在 $S$ 上的分离代数空间 $U$ 上固有作用, 且 $G\to S$ 是约化群概形, 或者 $S=\Spec k$, $\operatorname{char}(k)>0$, 则存在分离几何商 $U/G$~\cite[定理 $1.5$]{Kol97}. 在这些文章之前, 粗模空间的存在性被视为熟知定理, 但分离 Deligne--Mumford 叠情形的陈述和证明无疑早已为专家所知. 例如, \cite[定理 I.4.10]{FC90} 和~\cite[定理 $2.2$]{Kol90} 中明确陈述了此定理, 但未给出证明. Popp 在~\cite[第 $55$ 页的定理]{Pop74} 中, 对 $\mathbb C$ 上群作用的情形勾勒了使用解析技巧的论证, 而~\cite[\S$5$.A]{MF82} 中勾勒了解析和代数两种论证.

由于 $\mathcal M_g$ 是 $\mathbb Z$ 上的分离有限型 Deligne--Mumford 叠 (推论~\bookref{4.6.10} 和习题~\bookref{4.8.12}), Keel--Mori 定理蕴含存在作为分离代数空间的粗模空间 $\mathcal M_g\to M_g$. 此结果最初归功于 Mumford, 他证明了更强的结果, 即 $M_g$ 作为概形存在. 在~\cite{Mum61} 中勾勒证明之后, Mumford 在~\cite[定理 $5.11$ 和 $7.13$]{GIT} 中给出 $\Spec\mathbb Z$ 上 $M_g$ 的两种构造. 此外, 他证明对每个素数 $p$, $M_g$ 在 $\Spec\mathbb Z[1/p]$ 上拟射影. 后来证明了 $\overline M_g$ 在 $\Spec\mathbb Z$ 上射影, 特别地 $M_g$ 在 $\Spec\mathbb Z$ 上拟射影; 见~\S\bookref{6.8}--\bookref{6.9}. 有趣的是, Mumford 的两种构造实际都没有使用 GIT, 或至少没有使用通常被视为``标准 GIT 工具''的部分, 我们将在~\S\bookref{8.4} 中发展这些工具. Mumford 的一种构造依赖主极化 Abel 簇的模空间 $A_g$ 的存在性, 另一种则依赖使用协变式的专门方法.
\end{remark}
\begin{remark}\label{src:5.4.8}
强调此定理不仅给出粗模空间, 而且给出的粗模空间满足性质 (1)--(3). 仅满足定义中性质 (定义~\bookref{5.2.1}) 的粗模空间可能难以使用. 例如, 一般情况下尚不清楚粗模空间是否在平展基变换 (甚至开浸入) 下稳定, 或是否满足 $\pi_*\mathcal O_{\mathcal X}=\mathcal O_X$.
\end{remark}
\begin{proof}
先处理 $S=\Spec R$ 仿射的情形. 问题在 $\mathcal X$ 上是 Zariski 局部的: 若 $\{\mathcal X_i\}$ 是 $\mathcal X$ 的 Zariski 开覆盖, 且各项有粗模空间 $\mathcal X_i\to X_i$, 则由于粗模空间唯一 (定义~\bookref{5.2.1}(2)), 每个交 $\mathcal X_i\cap\mathcal X_j$ 都有粗模空间 $U_{ij}$, 可与 $X_i$ 和 $X_j$ 的开子概形等同, 从而可粘合得到代数空间 $X$ 以及映射 $\mathcal X\to X$. 由引理~\bookref{5.2.10}, 它是粗模空间. 因此只需说明每个闭点 $x\in|\mathcal X|$ 都有一个存在粗模空间的开邻域.

由 Deligne--Mumford 叠的局部结构定理 (\bookref{5.3.1}), 存在仿射平展态射
\[
f:(\mathcal W=[\Spec A/G_x],w)\longrightarrow(\mathcal X,x),
\]
使得 $f$ 在 $w$ 处诱导几何稳定子群的同构. 断言由满足 $f$ 在 $z$ 处诱导几何稳定子群同构的点 $z\in|\mathcal W|$ 组成的轨迹 $\mathcal U$ 是开的. 为证明这一点, 利用相对惯性 $I_{\mathcal X}\to\mathcal X$ 有限这一事实, 它由 $\mathcal X$ 的分离性推出 (因为对角线固有且拟有限). $\mathcal W$ 上相对群概形的自然态射 $I_{\mathcal W}\to I_{\mathcal X}\times_{\mathcal X}\mathcal W$ 在 $z\in\mathcal W(k)$ 上方的纤维正是稳定子的态射 $G_z\to G_{f(z)}$. 将使用习题~\bookref{4.2.15} 中的 Cartesian 图
\[
\begin{tikzcd}
I_{\mathcal W}\arrow[r,"\Psi"]\arrow[d]&I_{\mathcal X}\times_{\mathcal X}\mathcal W\arrow[d]\\
\mathcal W\arrow[r]\arrow[ur,phantom,"\square" description]&\mathcal W\times_{\mathcal X}\mathcal W.
\end{tikzcd}
\]
由于 $\mathcal W\to\mathcal X$ 可表出, 平展且分离, 对角线 $\mathcal W\to\mathcal W\times_{\mathcal X}\mathcal W$ 是开闭浸入, 因而 $\Psi$ 也是如此. 由于 $I_{\mathcal X}\to\mathcal X$ 有限, $p_2:I_{\mathcal X}\times_{\mathcal X}\mathcal W\to\mathcal W$ 也有限. 因此
\[
p_2\bigl(|I_{\mathcal X}\times_{\mathcal X}\mathcal W|\setminus|I_{\mathcal W}|\bigr)\subseteq|\mathcal W|
\]
是闭的, 它的补集与轨迹 $\mathcal U$ 等同, 因而是开的. 设 $\pi_{\mathcal W}:\mathcal W\to W=\Spec A^{G_x}$ 是粗模空间 (定理~\bookref{5.2.3}). 选取包含 $\pi_{\mathcal W}(w)$ 的仿射开子概形 $X_1\subseteq\pi(\mathcal U)$. 则 $\mathcal X_1=\pi_{\mathcal W}^{-1}(X_1)$ 同构于商叠 $[\Spec A_1/G_x]$, 且 $X_1=\Spec A_1^{G_x}$. 这给出仿射平展态射
\[
g:(\mathcal X_1=[\Spec A_1/G_x],w)\longrightarrow(\mathcal X,x),
\]
在所有几何稳定子群上诱导双射.

现在说明开子叠 $\mathcal X_0:=\operatorname{im}(f)$ 有粗模空间. 由于 $g$ 仿射且 $\mathcal X_1=[\Spec A_1/G_x]$, 可将 $\mathcal X_2:=\mathcal X_1\times_{\mathcal X}\mathcal X_1$ 表为 $[\Spec A_2/G_x]$, 并有粗模空间 $\pi_2:\mathcal X_2\to X_2=\Spec A_2^{G_x}$. 由粗模空间的泛性质, 有图
\begin{equation}\label{src:5.4.9}
\begin{tikzcd}[column sep=large]
\mathcal X_2\arrow[r,shift left]\arrow[r,shift right]\arrow[d,"\pi_2"']&\mathcal X_1\arrow[r,"g"]\arrow[d,"\pi_1"]&\mathcal X_0=\operatorname{im}(g)\arrow[d,dashed,"\pi_0"]\\
X_2\arrow[r,shift left]\arrow[r,shift right]\arrow[ur,phantom,"\square" description]&X_1\arrow[r,dashed]&X_0,
\end{tikzcd}
\end{equation}
其中自然方块交换. 由于 $g$ 在所有点的几何稳定子群上诱导双射, 每个投影 $\mathcal X_2\to\mathcal X_1$ 也如此. 推论~\bookref{5.4.5} 蕴含每个映射 $X_2\to X_1$ 平展, 且~\bookeqref{5.4.9} 中实线箭头的自然方块是 Cartesian 的.

粗模空间的泛性质诱导平展群胚结构 $X_2\rightrightarrows X_1$. 为检验它是平展等价关系, 只需检验 $X_2\to X_1\times X_1$ 在几何点上单射, 这由 $|\mathcal X_2|\to|\mathcal X_1|\times|\mathcal X_1|$ 在闭点上单射这一观察推出. 因此存在代数空间商 $X_0:=X_1/X_2$ 以及映射 $X_1\to X_0$. 沿 $\mathcal X_1\to\mathcal X_0$ 作平展下降, 得到映射 $\pi_0:\mathcal X_0\to X_0$, 使~\bookeqref{5.4.9} 的右方块交换.

为说明 $\pi:\mathcal X_0\to X_0$ 是粗模空间, 将使用交换立方体
\[
\begin{tikzcd}[row sep=1.4em,column sep=1.6em]
&\mathcal X_2\arrow[rr]\arrow[dl]\arrow[dd]&&\mathcal X_1\arrow[dl]\arrow[dd]\\
\mathcal X_1\arrow[rr,crossing over]\arrow[dd]&&\mathcal X_0\arrow[dd]&\\
&X_2\arrow[rr]\arrow[dl]&&X_1\arrow[dl]\\
X_1\arrow[rr]&&X_0.&
\end{tikzcd}
\]
其中上面, 左面和下面是 Cartesian 的. 由于 $\mathcal X_1\to\mathcal X_0$ 平展且满射, 自然映射 $\mathcal X_1\to\mathcal X_0\times_{X_0}X_1$ 拉回后为同构, 故 $\mathcal X_1\cong\mathcal X_0\times_{X_0}X_1$, 右方块也为 Cartesian 的. 由于为粗模空间的性质在 $X_0$ 上是平展局部的 (引理~\bookref{5.2.10}), 得到 $\pi_0:\mathcal X_0\to X_0$ 是粗模空间. $\mathcal O_{X_0}\to\pi_{0,*}\mathcal O_{\mathcal X_0}$ 的条件及性质 (1)--(3) 也都是平展局部的, 且由定理~\bookref{5.2.3}, 它们对 $[\Spec A_1/G_x]\to\Spec A_1^{G_x}$ 成立, 故也对 $\mathcal X_0\to X_0$ 成立.

$S$ 为 Noether 代数空间的情形, 可仿照上述论证, 在平展局部构造 $\mathcal X$ 的粗模空间, 从而归结到仿射情形. 最后, 由于 $\mathcal X_0\to X_0$ 固有, $\mathcal X_0$ 的分离性等价于 $X_0$ 的分离性.
\end{proof}
\begin{exercise}\label{src:5.4.10}
若 $\mathcal X$ 是 Noether 代数空间 $S$ 上的有限型 Deligne--Mumford 叠, 且 $I_{\mathcal X}\to\mathcal X$ 有限, 则存在粗模空间 $\pi:\mathcal X\to X$, 满足 $\mathcal O_X=\pi_*\mathcal O_{\mathcal X}$, 并具有定理~\bookref{5.4.6} 的性质 (1)--(3).
\end{exercise}
\begin{remark}\label{src:5.4.11}
$\mathcal X$ 为具有有限惯性 $I_{\mathcal X}\to\mathcal X$ 的代数叠 (不一定为 Deligne--Mumford 叠) 的更一般情形, 可用类似但更为技术性的方式证明:
\begin{enumerate}[{步骤} 1:]
\item 构造来自概形的拟有限平坦满射 $U\to\mathcal X$. 证明与极小呈示的存在性 (\bookref{4.6.1}) 类似; 见注记~\bookref{4.6.2}, \cite[引理 V.7.2]{SGA3I}, \cite[引理 $3.3$]{KM97}, \cite[引理 $2.1$]{Con05b} 和~\cite[Tag 06MC]{SP}.
\item 说明每个点都有平展邻域 $\mathcal W\to\mathcal X$, 其中 $\mathcal W$ 有来自仿射概形的有限平坦呈示 $V\to\mathcal W$. 这应当视为 Deligne--Mumford 叠局部结构定理 (\bookref{5.3.1}) 的替代: $\mathcal W$ 不再是有限群的商 $[\Spec A/G]$, 而是仿射概形的有限平坦群胚 $R:=V\times_{\mathcal W}V\rightrightarrows V$ 的商 $[V/R]$. 证明也类似: 不使用 $\mathrm{ET}_d:=[(U\times_{\mathcal X}\cdots\times_{\mathcal X}U\setminus\Delta)/S_d]$ 参数化 $U\to\mathcal X$ 的纤维中的 $d$ 个无序点, 而使用 $\mathcal X$ 上的相对 Hilbert 叠 $\mathcal{H}ilb_d$, 其在 $S\to\mathcal X$ 上的对象由闭子概形 $Z\hookrightarrow U\times_{\mathcal X}S$ 组成, 在 $S$ 上有限平坦 (而非有限平展), 次数为 $d$. 见~\cite[\S$5$.A]{MF82}, \cite[命题 $4.2$]{KM97}, \cite[引理 $2.2$]{Con05b} 和~\cite[Tag 0DU4]{SP}.
\item 证明仿射概形的有限平坦群胚 $R\rightrightarrows V$ 的商 $[R/V]$ 的群胚存在性. 证明与定理~\bookref{5.2.3} 类似; 见习题~\bookref{5.2.16}, \cite[定理 V.4.1]{SGA3I}, \cite[命题 $5.1$]{KM97}, \cite[定理 $3.1$]{Con05b}, \cite[定理 $6.2.2$]{Ols16} 和~\cite[Tag 0DUR]{SP}.
\end{enumerate}
\end{remark}
\begin{corollary}\label{src:5.4.12}
若 $G$ 是 Noether 概形 $S$ 上的仿射 fppf 群概形, 在分离代数空间 $U$ 上固有作用, 即作用映射 $G\times U\to U\times U$ 固有, 且所有稳定子群均约化, 则存在到 $S$ 上分离有限型代数空间 $X$ 的几何商 $\pi:U\to X$, 满足 $\mathcal O_X=(\pi_*\mathcal O_U)^G$.
\end{corollary}
\begin{proof}
由于 $G$ 在 $S$ 上仿射, $\sigma:G\times U\to U\times U$ 仿射. 因此 $\sigma$ 的固有性等价于 $U$ 的有限性, 并蕴含 $[U/G]$ 的对角线有限. 由于稳定子群约化, $[U/G]$ 是分离的 Deligne--Mumford 叠, 从而可应用 Keel--Mori 定理 (\bookref{5.4.6}).
\end{proof}
证明中粗模空间的构造给出粗模空间的局部结构定理. (也可直接由局部结构定理 (\bookref{5.3.1}) 和习题~\bookref{5.4.14} 得到以下定理.)
\begin{corollary}[粗模空间的局部结构]\label{src:5.4.13}
设 $\mathcal X$ 是 Noether 代数空间 $S$ 上的有限型分离 Deligne--Mumford 叠, $\pi:\mathcal X\to X$ 是其粗模空间. 对每个闭点 $x\in|\mathcal X|$, 设其几何稳定子群为 $G_x$, 则存在 Cartesian 图
\[
\begin{tikzcd}
{[\Spec A/G_x]}\arrow[r]\arrow[d]&\mathcal X\arrow[d,"\pi"]\\
\Spec A^{G_x}\arrow[r]\arrow[ur,phantom,"\square" description]&X,
\end{tikzcd}
\]
使得 $\Spec A^{G_x}\to X$ 是 $\pi(x)\in|X|$ 的平展邻域.
\end{corollary}
\begin{exercise}\label{src:5.4.14}
建立命题~\bookref{5.4.1} 的以下推广: 若 $f:\mathcal X\to\mathcal Y$ 是 Noether 代数空间 $S$ 上分离有限型 Deligne--Mumford 叠的态射, 位于交换图
\[
\begin{tikzcd}
\mathcal X\arrow[r,"f"]\arrow[d,"\pi_{\mathcal X}"']&\mathcal Y\arrow[d,"\pi_{\mathcal Y}"]\\
X\arrow[r]&Y
\end{tikzcd}
\]
中, 其中 $\pi_{\mathcal X}:\mathcal X\to X$ 和 $\pi_{\mathcal Y}:\mathcal Y\to Y$ 是粗模空间, 且 $x\in|\mathcal X|$ 是满足以下条件的闭点:
\begin{enumerate}[(1)]
\item $f$ 在 $x$ 处平展; 且
\item 所诱导的几何稳定子群映射 $G_x\to G_{f(x)}$ 是双射,
\end{enumerate}
则 $\pi_{\mathcal X}(x)$ 有开邻域 $U\subseteq X$, 使得 $U\to X\to Y$ 平展, 且 $\pi_{\mathcal X}(U)\cong U\times_Y\mathcal Y$.
\end{exercise}

\subsection{例子}\label{src:5.4.3-subsection}
粗模空间的基本例子是有限群作用所诱导的映射 $[\Spec A/G]\to\Spec A^G$; 一些初等例子见例~\bookref{5.2.4}.
\begin{example}[椭圆曲线]\label{src:5.4.15}
考虑习题~\bookref{4.1.21} 中椭圆曲线的模叠 $\mathcal M_{1,1}\cong[(\mathbb A^2\setminus V(\Delta))/\mathbb G_m]$, 其中 $k$ 是满足 $\operatorname{char}(k)\neq2,3$ 的域, $\mathbb A^2$ 的坐标为 $a,b$, 作用为 $t\cdot(a,b)=(t^4a,t^6b)$, $\Delta=4a^3+27b^2$ 是判别式. 不变量环 $k[a,b]_\Delta^{\mathbb G_m}$ 是 $k[a^3/\Delta]$, 其中 $\beta:=b^2/\Delta$ 由 $\alpha:=a^3/\Delta$ 通过关系 $4\alpha+27\beta=1$ 生成. 因此粗模空间 $\mathcal M_{1,1}\to\mathbb A^1$ 由 $(a,b)\mapsto a^3/\Delta$ 给出. 对 $\lambda\in k$, 由 $\mathbb P^2$ 中 Legendre 方程 $y^2z=x(x-z)(x-\lambda z)$ 表示的椭圆曲线 $E_\lambda$, 射影等价于某个 $a_\lambda,b_\lambda\in k$ 对应的方程 $y^2z=x^3+a_\lambda xz^2+b_\lambda z^3$, 且可检验 $a_\lambda^3/\Delta_\lambda$ 是 $j$-不变量
\[
2^8\frac{(\lambda^2-\lambda+1)^3}{\lambda^2(\lambda-1)^3}
\]
的标量倍数. 也可使用习题~\bookref{4.4.21} 中的商呈示 $\mathcal M_{1,1}\cong[V(\Delta-1)/\mu_{12}]$, 粗模空间同样由 $(a,b)\mapsto a^3/\Delta$ 给出.
\end{example}
\begin{exercise}[困难]\label{src:5.4.16}
设 $\operatorname{char}(k)\neq2$, $G=\mathbb Z/2$.
\begin{enumerate}[(a)]
\item 令 $G$ 通过交换两个 $\mathbb A^1$ 副本作用于非分离的并 $X=\mathbb A^1\bigcup_{x\neq0}\mathbb A^1$. 商 $[X/G]$ 是具有拟有限但非有限惯性的 Deligne--Mumford 叠, 特别地它不分离. 证明仍存在粗模空间 $[X/G]\to\mathbb A^1$.
\item 设 $X$ 是非分离的并 $\mathbb A^2\bigcup_{x\neq0}\mathbb A^2$. 令 $G=\mathbb Z/2$ 同时交换两个 $\mathbb A^2$ 副本, 并在每个副本上通过对合 $y\mapsto-y$ 作用于 $X$. 证明 $[X/G]$ 没有粗模空间.
\end{enumerate}
\end{exercise}
\newpage
\begin{example}\label{src:5.4.17}
考虑 $\PGL_2$ 在概形 $\Sym^4\mathbb P^1\cong(\mathbb P^1)^4/S_4$ 上的作用, 该概形参数化 $\mathbb P^1$ 中的四个无序点, 并且是 $[(\mathbb P^1)^4/S_4]$ 的粗模空间. 设 $\mathcal X\subseteq[\Sym^4\mathbb P^1/\PGL_2]$ 是参数化至少有三个点不同的元组 $(p_1,p_2,p_3,p_4)$ 的开子叠. 考虑 $\lambda\in\mathbb P^1$ 给出的族 $(0,1,\lambda,\infty)$. 若 $\lambda\notin\{0,1\infty\}$, 则断言 $\operatorname{Aut}(0,1,\lambda,\infty)=\mathbb Z/2\times\mathbb Z/2$. 为说明这一点, 有唯一元素 $\sigma\in\PGL_2$, 满足 $\sigma(0)=\infty$, $\sigma(\infty)=0$ 和 $\sigma(1)=\lambda$, 它通过 $\sigma([x,y])=[y,\lambda,x]$ 作用于 $\mathbb P^1$, 因而 $\sigma(\lambda)=1$. 类似地, 存在交换 $0$ 与 $1$, 同时交换 $\lambda$ 与 $\infty$ 的元素, 也存在交换 $0$ 与 $\lambda$, 同时交换 $1$ 与 $\infty$ 的元素. 然而, 若 $\lambda\in\{0,1\infty\}$, 则 $\operatorname{Aut}(0,1,\lambda,\infty)=\mathbb Z/2$. 因此惯性 $I_{\mathcal X}\to\mathcal X$ 虽然拟有限但并不有限, $\mathcal X$ 不分离. 尽管如此, 将 $(p_1,p_2,p_3,p_4)$ 送到其交比的映射 $\mathcal X\to\mathbb P^1$ 是粗模空间.
\end{example}
\begin{exercise}\label{src:5.4.18}
设 $\pi:\mathcal X\to X$ 是 Noether 仿射概形上分离有限型 Deligne--Mumford 叠 $\mathcal X$ 的粗模空间.
\end{exercise}
\begin{exercise}\label{src:5.4.19}
设 $\mathcal X$ 是 Noether 代数空间 $S$ 上的有限型分离 Deligne--Mumford 叠, $\pi:\mathcal X\to X$ 是其粗模空间.
\begin{enumerate}[(a)]
\item 若 $x\in|\mathcal X|$ 是有限型点, 证明 $x$ 的剩余 gerbe $\mathcal G_x$ 可与纤维 $\mathcal X_{\pi(x)}$ 的约化等同, 并给出纤维 $\mathcal X_{\pi(x)}$ 非约化的例子.
\item 证明若 $\mathcal X$ 正规, 则 $X$ 也正规.
\item 若 $\mathcal X$ 正则, 证明 $\mathcal X\to X$ 平坦, 当且仅当 $X$ 正则.
\item 给出粗模空间 $\mathcal X\to X$ 不平坦的例子.
\end{enumerate}
\end{exercise}

\subsection{温和粗模空间}\label{src:5.4.4-subsection}
\begin{definition}\label{src:5.4.20}
Deligne--Mumford 叠 $\mathcal X$ 称为\newterm{温和}{tame}, 若对每个几何点 $x\in\mathcal X(k)$, $\operatorname{Aut}_{\mathcal X(k)}(x)$ 的阶在 $\Gamma(\mathcal X,\mathcal O_{\mathcal X})$ 中可逆. 若此外 $\mathcal X$ 有粗模空间 $\mathcal X\to X$, 则称 $\mathcal X\to X$ 为\newterm{温和粗模空间}{tame coarse moduli space}.
\end{definition}
\begin{remark}\label{src:5.4.21}
若 $\mathcal X$ 定义在域 $k$ 上, 则这意味着每个几何稳定子群的阶都与 $k$ 的特征互素.
\end{remark}
\begin{lemma}\label{src:5.4.22}
设 $\mathcal X$ 是 Noether 代数空间 $S$ 上的分离有限型 Deligne--Mumford 叠. 若 $\pi:\mathcal X\to X$ 是温和粗模空间, 则 $\pi_*$ 正合.
\end{lemma}
\begin{proof}
问题在 $X$ 上是平展局部的: 若 $g:X'\to X$ 是平展覆盖, 诱导 Cartesian 图
\[
\begin{tikzcd}
\mathcal X'\arrow[r,"{g'}"]\arrow[d,"{\pi'}"']&\mathcal X\arrow[d,"\pi"]\\
X'\arrow[r,"g"']\arrow[ur,phantom,"\square" description]&X,
\end{tikzcd}
\]
则由平坦基变换 (\bookref{5.1.43}), 拟凝聚层上的函子满足等同 $g^*\pi_*=\pi'_*g'^*$. 由于 $g^*$ 忠实正合, $\pi_*$ 正合, 当且仅当 $\pi'_*$ 正合. 因而可用局部结构定理 (\bookref{5.4.13}) 归结到 $\mathcal X=[\Spec A/G]$, $X=\Spec A^G$ 的情形. 此时 $\pi_*$ 的正合性由它可解释为取不变量 (习题~\bookref{5.1.27}) 以及 $G$ 的线性约化性, 即 Maschke 定理 (\bookref{B.1.41}), 推出.
\end{proof}
\begin{corollary}\label{src:5.4.23}
设 $\mathcal X$ 是 Noether 代数空间 $S$ 上的分离有限型 Deligne--Mumford 叠, $\pi:\mathcal X\to X$ 是其粗模空间. 若 $F$ 是 $\mathcal X$ 上的拟凝聚层, 则 $\mathrm H^i(\mathcal X,F)=\mathrm H^i(X,F)$, 且当 $i>\dim\mathcal X$ 时为 $0$.
\end{corollary}
\begin{proof}
第一个陈述由 $\pi_*$ 的正合性推出, 第二个由习题~\bookref{5.1.48} 推出.
\end{proof}
\begin{exercise}[基变换]\label{src:5.4.24}
设 $\mathcal X$ 是 Noether 代数空间 $S$ 上的有限型分离 Deligne--Mumford 叠. 若 $\pi:\mathcal X\to X$ 是温和粗模空间, $X'\to X$ 是代数空间的态射, 证明 $\mathcal X\times_X X'\to X'$ 也是粗模空间.
\end{exercise}
另一方面, 习题~\bookref{5.2.8}(c) 给出一个粗模空间 $\mathcal X\to X$ 和映射 $X'\to X$, 使得 $\mathcal X\times_X X'\to X'$ 不是粗模空间. 特别地, 不能自动推出 $\overline M_g\times_{\mathbb Z}\mathbb F_p$ 是 $\overline{\mathcal M}_g\times_{\mathbb Z}\mathbb F_p$ 的粗模空间; 见问题~\bookref{6.5.25}.
\begin{remark}[典范叠与粗模空间的自下而上分类]\label{src:5.4.25}
域 $k$ 上光滑分离温和 Deligne--Mumford 叠 $\mathcal X$ 的粗模空间 $X$ 具有\newterm{温和商奇点}{tame quotient singularities}, 即 $X$ 的每个点都有平展邻域, 它是阶与 $\operatorname{char}(k)$ 互素的有限群作用于光滑概形 $V$ 的几何商 $V/G$. 这由粗模空间的局部结构 (\bookref{5.4.13}) 推出. 反之, Vistoli 证明, 对每个具有温和商叠的代数空间 $X$, 存在以 $X$ 为粗模空间的典范光滑分离温和 Deligne--Mumford 叠 $\mathcal X^{\mathrm{can}}$~\cite[命题 $2.8$]{Vis89}. 对任意具有温和商奇点的分离 Deligne--Mumford 叠, 也存在相对\newterm{典范叠}{canonical stack} $\mathcal X^{\mathrm{can}}\to\mathcal X$. Geraschenko 和 Satriano 的一个定理表明, 每个一般稳定子平凡的光滑分离温和 Deligne--Mumford 叠 $\mathcal X$, 都可从其粗模空间 $X$ 恢复为复合
\[
\mathcal X\cong\bigl(X^{\mathrm{can}}(\sqrt D)\bigr)^{\mathrm{can}}\longrightarrow X^{\mathrm{can}}(\sqrt D)\longrightarrow X^{\mathrm{can}}\longrightarrow X,
\]
其中 $X^{\mathrm{can}}(\sqrt D)\to X^{\mathrm{can}}$ 是沿 $\mathcal X\to X$ 的分歧除子的不可约分支所取的根叠 (例~\bookref{4.9.22}) 的复合, 每次取根的阶为分歧次数~\cite{GS15}.
\end{remark}

\subsection{将向量丛下降到粗模空间}\label{src:5.4.5-subsection}
先给出凝聚层的一个 Nakayama 引理.
\begin{lemma}\label{src:5.4.26}
设 $\mathcal X$ 是 Noether 代数空间 $S$ 上的分离有限型 Deligne--Mumford 叠, $\pi:\mathcal X\to X$ 是其粗模空间. 设 $x\in|\mathcal X|$ 是闭点, $\mathcal G_x$ 是其剩余 gerbe.
\begin{enumerate}[(1)]
\item 若 $F$ 是 $\mathcal X$ 上满足 $F|_{\mathcal G_x}=0$ 的凝聚层, 则 $\pi(x)$ 有开邻域 $U\subseteq X$, 使得 $F|_{\pi^{-1}(U)}=0$.
\item 若 $\phi:F\to G$ 是 $\mathcal X$ 上凝聚层 (分别为同秩向量丛) 的态射, 且 $\phi|_{\mathcal G_x}$ 满射, 则 $\pi(x)$ 有开邻域 $U\subseteq X$, 使得 $\phi|_{\pi^{-1}(U)}$ 满射 (分别为同构).
\end{enumerate}
\end{lemma}
\begin{proof}
对 (1), $F$ 的支集 $\operatorname{Supp}(F)\subseteq|\mathcal X|$ 是闭子集 (这由沿呈示的下降推出), 开集 $U=X\setminus\pi(\operatorname{Supp}(F))$ 满足结论. 对 (2), 将 (1) 应用于凝聚层 $\operatorname{coker}(\phi)$, 并注意同秩向量丛的满射是同构.
\end{proof}
称 $\mathcal X$ 上的向量丛 $F$ 下降到其粗模空间 $\pi:\mathcal X\to X$, 若存在 $X$ 上的向量丛 $G$ 以及同构 $F\cong\pi^*G$. 注意一个必要条件是: 对每个几何点 $x:\Spec k\to\mathcal X$ 以及交换图
\[
\begin{tikzcd}
BG_x\arrow[r,hook,"i_x"]\arrow[d,"p"']&\mathcal X\arrow[d,"\pi"]\\
\Spec k\arrow[r,hook,"j"']&X,
\end{tikzcd}
\]
拉回 $i_x^*F=i_x^*\pi^*G=p^*j^*G$ 平凡, 换言之 $G_x$ 在纤维 $F\otimes k:=i_x^*F$ 上作用平凡. 这等价于要求 $F$ 在 $x$ 的剩余 gerbe 上的限制 $F|_{\mathcal G_x}$ 平凡.
\begin{proposition}\label{src:5.4.27}
设 $\mathcal X$ 是 Noether 代数空间 $S$ 上的温和分离有限型 Deligne--Mumford 叠, $\pi:\mathcal X\to X$ 是其粗模空间. $\mathcal X$ 上的向量丛 $F$ 下降到 $X$ 上的向量丛, 当且仅当对每个具有闭像的几何点 $x:\Spec k\to\mathcal X$, $G_x$ 在纤维 $F\otimes k$ 上作用平凡. 此时 $\pi_*F$ 是向量丛, 伴随映射 $\pi^*\pi_*F\to F$ 是同构.
\end{proposition}
\begin{proof}
为说明条件充分, 考虑交换图
\[
\begin{tikzcd}
\mathcal G_x\arrow[r,hook]\arrow[d,"p"']&\mathcal X\arrow[d,"\pi"]\\
\Spec\kappa(x)\arrow[r,hook]&X,
\end{tikzcd}
\]
其中 $\mathcal G_x$ 是剩余 gerbe. 将证明分为三个步骤.
\begin{enumerate}[{步骤} 1:]
\item $\pi^*\pi_*F\to F$ 满射. 由引理~\bookref{5.4.26}, 只需对每个闭点 $x\in|\mathcal X|$ 说明 $(\pi^*\pi_*F)|_{\mathcal G_x}\to F|_{\mathcal G_x}$ 满射. 由于 $F\to F|_{\mathcal G_x}$ 满射且 $\pi_*$ 正合 (引理~\bookref{5.4.22}), $\pi_*F\to\pi_*(F|_{\mathcal G_x})$ 满射, 因而
\[
(\pi^*\pi_*F)|_{\mathcal G_x}\longrightarrow\pi^*\bigl(\pi_*(F|_{\mathcal G_x})\bigr)|_{\mathcal G_x}\cong p^*p_*(F|_{\mathcal G_x})
\]
也满射. 假设蕴含伴随映射 $p^*p_*(F|_{\mathcal G_x})\to F|_{\mathcal G_x}$ 为同构, 从而 $(\pi^*\pi_*F)|_{\mathcal G_x}\to p^*p_*(F|_{\mathcal G_x})\cong F|_{\mathcal G_x}$ 满射.
\item $\pi_*F$ 是向量丛. 可假设 $F$ 的秩 $r$ 为常数. 由于为向量丛的性质是平展局部的, 可假设 $X=\Spec A$. 满射 $\bigoplus_{s\in\Gamma(X,\pi_*F)}A\to\pi_*F$ 拉回为满射 $\bigoplus_{s\in\Gamma(\mathcal X,F)}\mathcal O_{\mathcal X}\to\pi^*\pi_*F$. 由步骤 1, 复合 $\bigoplus_{s\in\Gamma(\mathcal X,F)}\mathcal O_{\mathcal X}\to\pi^*\pi_*F\to F$ 满射. 由于 $F|_{\mathcal G_x}\cong\mathcal O_{\mathcal G_x}^r$ 平凡, 对每个闭点 $x\in|\mathcal X|$, 可找到 $r$ 个截面 $\phi:\mathcal O_{\mathcal X}^{\oplus r}\to F$, 使得 $\phi|_{\mathcal G_x}$ 是同构. 由引理~\bookref{5.4.26}, $\pi(x)$ 有开邻域 $U\subseteq X$, 使得 $\phi|_{\pi^{-1}(U)}$ 是同构. 因此 $\pi_*\phi:\mathcal O_X^{\oplus r}\to\pi_*F$ 在 $U$ 上是同构, 从而 $\pi_*F$ 是与 $F$ 同秩的向量丛.
\item $\pi^*\pi_*F\to F$ 是同构. 由于 $\pi^*\pi_*F\to F$ 是同秩向量丛的满射, 它是同构.
\end{enumerate}
\end{proof}
\begin{remark}\label{src:5.4.28}
对凝聚层, 相应陈述并不成立. 例如, 在特征不为 $2$ 时, 令 $\mathbb Z/2$ 通过 $x\mapsto-x$ 作用于 $\mathbb A^1$, 有温和粗模空间 $[\mathbb A^1/(\mathbb Z/2)]\to\mathbb A^1=\Spec k[x^2]$. 原点 $0$ 的包含 $B(\mathbb Z/2)\hookrightarrow[\mathbb A^1/(\mathbb Z/2)]$ 是闭子叠, $\mathcal O_{B(\mathbb Z/2)}$ 是不下降的凝聚层.
\end{remark}
当 $\mathcal X$ 不温和时, 线丛下降有以下变体.
\begin{exercise}\label{src:5.4.29}
设 $\mathcal X$ 是 Noether 代数空间 $S$ 上的分离有限型 Deligne--Mumford 叠, $\pi:\mathcal X\to X$ 是其粗模空间. 设 $L$ 是 $\mathcal X$ 上的线丛.
\begin{enumerate}[(1)]
\item 证明对充分可除的 $N$, $L^{\otimes N}$ 下降到 $X$.

\textit{提示:} 证明对充分可除的 $N$, $\pi_*(L^{\otimes N})$ 是线丛, 且伴随映射 $\pi^*\pi_*(L^{\otimes N})\to L^{\otimes N}$ 是同构. 归结到 $\mathcal X=[U/G]$ 且 $L$ 拉回到 $U$ 上为平凡线丛的情形, 并使用由复合 $[U/G]\to BG\to\Spec\mathbb Z$ 诱导的谱序列
\[
E_2^{pq}=\mathrm H^p\bigl(BG,\mathrm H^q(U,\mathbb G_m)\bigr)\Longrightarrow\mathrm H^{p+q}([U/G],\mathbb G_m).
\]
\item 假设 $\mathcal X$ 在 $S$ 上的有限型 Deligne--Mumford 叠 $\mathcal Y$ 上固有且可表出. 若 $L$ 在 $\mathcal Y$ 上相对丰沛, 证明对充分可除的 $N$, $L^{\otimes N}$ 下降为在 $\mathcal Y$ 的粗模空间 $Y$ 上相对丰沛的线丛.

(丰沛性是底空间上平展局部的性质. 称 $\mathcal X$ 上的线丛 $L$ 在 $\mathcal Y$ 上相对丰沛, 若 $\mathcal X\to\mathcal Y$ 可由概形表出, 且存在平展呈示 $V\to\mathcal Y$, 使得 $L$ 在 $\mathcal X_V$ 上的限制 $L_V$ 相对丰沛.)
\end{enumerate}
\end{exercise}
